\documentclass[11pt,openright]{book}

\usepackage[paperwidth=6.5in,paperheight=9.25in,top=0.78in,bottom=0.72in,inner=0.82in,outer=0.68in,headheight=23pt]{geometry}
\usepackage[T1]{fontenc}
\usepackage[utf8]{inputenc}
\usepackage{lmodern}
\usepackage{microtype}
\usepackage{amsmath,amssymb}
\usepackage{booktabs,tabularx,longtable,array}
\usepackage{xcolor}
\usepackage{graphicx}
\usepackage{tikz}
\usetikzlibrary{arrows.meta,positioning,shapes.geometric,calc}
\usepackage{listings}
\usepackage{enumitem}
\usepackage{url}
\usepackage{makeidx}
\usepackage{fancyhdr}
\usepackage{hyperref}

\definecolor{ink}{HTML}{17212B}
\definecolor{muted}{HTML}{53616D}
\definecolor{accent}{HTML}{0B7285}
\definecolor{codebg}{HTML}{F3F6F8}
\definecolor{codekw}{HTML}{075985}
\definecolor{codecomment}{HTML}{4D7C0F}
\definecolor{rulegray}{HTML}{C9D2D9}

\hypersetup{
  colorlinks=true,
  linkcolor=accent,
  citecolor=accent,
  urlcolor=accent,
  pdftitle={Engineering Software},
  pdfauthor={Scott Brodie Forsyth}
}

\lstset{
  basicstyle=\ttfamily\small,
  backgroundcolor=\color{codebg},
  frame=single,
  rulecolor=\color{rulegray},
  framerule=0.4pt,
  framesep=5pt,
  xleftmargin=7pt,
  xrightmargin=7pt,
  breaklines=true,
  columns=fullflexible,
  keepspaces=true,
  showstringspaces=false,
  keywordstyle=\color{codekw}\bfseries,
  commentstyle=\color{codecomment}\itshape,
  numbers=left,
  numberstyle=\tiny\color{muted},
  numbersep=6pt
}

\pagestyle{fancy}
\fancyhf{}
\fancyhead[LE]{\small\sffamily\color{muted}\leftmark}
\fancyhead[RO]{\small\sffamily\color{muted}\rightmark}
\fancyfoot[C]{\small\sffamily\thepage}
\setlength{\headwidth}{\textwidth}

\setlist{itemsep=2pt,topsep=4pt}
\setlength{\tabcolsep}{3pt}
\setcounter{tocdepth}{1}
\setcounter{secnumdepth}{2}
\makeindex

\makeatletter
\renewcommand{\cleardoublepage}{%
  \clearpage
  \if@twoside
    \ifodd\c@page
    \else
      \hbox{}\thispagestyle{empty}\newpage
      \if@twocolumn\hbox{}\newpage\fi
    \fi
  \fi
}
\makeatother

\newcommand{\term}[1]{\textbf{#1}\index{#1}}
\newcommand{\artifact}[1]{\texttt{#1}}
\newenvironment{keyidea}{\begin{quote}\small\itshape\color{accent}}{\end{quote}}
\newenvironment{tradeoffs}{\begin{quote}\small\color{muted}\textbf{Trade-offs.}\ }{\end{quote}}
\newenvironment{exercise}{\begin{quote}\small\color{ink}\textbf{Exercise.}\ }{\end{quote}}

\begin{document}

\frontmatter
\begin{titlepage}
  \pagecolor{ink}
  \color{white}
  \vspace*{1.15in}
  {\sffamily\fontsize{24}{28}\selectfont\bfseries ENGINEERING SOFTWARE\par}
  \vspace{0.22in}
  {\sffamily\Large From Reliable Code to Production Systems\par}
  \vspace{0.36in}
  \textcolor{white!55}{\rule{1.2in}{0.8pt}}\par
  \vfill
  {\sffamily\large Scott Brodie Forsyth\par}
  \vspace{0.25in}
  {\sffamily\small A compact, evidence-aware guide to designing, building,\par\noindent\hspace*{0pt}testing, securing, deploying, and operating software.\par}
  \thispagestyle{empty}
\end{titlepage}
\nopagecolor

\cleardoublepage
\thispagestyle{empty}
\vspace*{1.2in}
\begin{center}
  {\Large\sffamily\bfseries Engineering software is the practice of making\par}
  \vspace{0.15in}
  {\large\sffamily\color{muted}important behaviour dependable under changing conditions.\par}
\end{center}
\vfill
\noindent\small
Copyright \textcopyright\ 2026 Scott Brodie Forsyth.\par
\vspace{0.1in}
This edition is prepared as an editable educational manuscript. Standards,
dependencies, and platform interfaces change; verify operational details before
using them in a live system.\par
\vspace{0.1in}
Source code examples are provided for learning and are licensed for reuse with
the repository. Third-party standards and books remain under their original
licenses.

\cleardoublepage
\tableofcontents

\chapter*{How to use this book}
\addcontentsline{toc}{chapter}{How to use this book}
Software engineering is not a bag of rituals. It is a way of reducing the gap
between what a system is supposed to do and what it actually does when people,
networks, data, deadlines, attackers, and future maintainers interact with it.

Each chapter moves between four levels: a mental model, a concrete technique,
the conditions under which the technique helps, and the way it can fail. The
running project is a small order-intake service called \emph{Ledgerline}. It is
intentionally ordinary: it accepts orders, validates them, stores them, exposes
an API, emits useful signals, and must remain understandable six months later.
The repository contains smaller examples so that the ideas can be run without
proprietary infrastructure.

Words such as \emph{must}, \emph{should}, and \emph{may} are used carefully.
When they are capitalized inside a quoted standard, they retain the meaning of
that standard. Elsewhere, they indicate the strength of the author's advice,
not a universal law.

\mainmatter
\chapter{Engineering as a way of reasoning}

\section{The object of the discipline}
Programming is the act of expressing a computation. \term{Software engineering}
adds a harder question: how can that computation remain correct, understandable,
secure, and economically maintainable while its environment changes? The
environment includes users, regulations, dependencies, hardware, networks,
other teams, and the people who will inherit the code without the original
author's memory.

The word \emph{engineering} is useful because it points to constrained design.
An engineer makes assumptions explicit, models important failure modes, chooses
materials and boundaries, measures the result, and manages uncertainty. Software
has unusual properties: it is easy to copy, cheap to modify in one place, and
often difficult to observe as a whole. A change that looks local can alter
security, latency, data interpretation, or the behaviour of a distant client.

\begin{keyidea}
Quality is not a single attribute. It is the degree to which a system satisfies
the requirements that matter in its context, including requirements about
change, operation, safety, privacy, and future understanding.
\end{keyidea}

\section{Claims, evidence, and recommendations}
Technical writing often mixes three kinds of statement. A specification defines
what an implementation must do. Empirical research reports observations under
specified conditions. Professional practice is accumulated experience that may
be useful without being a universal law. A recommendation is a judgment about a
particular context. Keeping these categories separate prevents cargo cults.

For example, HTTP semantics are defined by an IETF standard; the assertion that
an HTTP method is safe or idempotent is a protocol claim that can be checked
against that standard \cite{ietf2022http}. The belief that small pull requests
are easier to review is a practice hypothesis: often plausible, but dependent
on change shape, reviewer skill, and system complexity. The decision to use a
modular monolith for a small product is an engineering recommendation whose
strength comes from the constraints, not from the label.

\section{A system boundary is a decision}
Every system description has a boundary. Inside it are behaviours the team can
change directly; outside it are actors, services, devices, laws, and physical
conditions. A boundary is never purely technical. It encodes ownership and
responsibility.

The Ledgerline boundary is deliberately modest:

\begin{itemize}
  \item a caller submits an order containing a customer identifier and line items;
  \item the service validates domain invariants and returns a stable order identifier;
  \item a relational store records the accepted state;
  \item operators can tell whether the service is alive, useful, and overloaded;
  \item a failure does not silently create a second order when the caller retries.
\end{itemize}

The last point is not an implementation detail. It is a business property that
leads to an idempotency key, a database constraint, a response contract, and a
recovery test. Engineering begins when a vague concern becomes a checkable
property.

\section{The feedback loop}
Good development is a loop rather than a waterfall of documents:

\begin{center}
\begin{tikzpicture}[node distance=6mm,>=Latex,font=\small]
  \node[draw,rounded corners,fill=codebg,minimum width=22mm,minimum height=8mm] (need) {need};
  \node[draw,rounded corners,fill=codebg,minimum width=22mm,minimum height=8mm,right=of need] (model) {model};
  \node[draw,rounded corners,fill=codebg,minimum width=22mm,minimum height=8mm,right=of model] (build) {build};
  \node[draw,rounded corners,fill=codebg,minimum width=22mm,minimum height=8mm,right=of build] (observe) {observe};
  \draw[->,thick,accent] (need) -- (model);
  \draw[->,thick,accent] (model) -- (build);
  \draw[->,thick,accent] (build) -- (observe);
  \draw[->,thick,accent] (observe) to[bend left=35] (need);
\end{tikzpicture}
\end{center}

The model may be a test, a state machine, a schema, a sequence diagram, or a
small written invariant. The point is not to predict every detail. It is to
make the important assumptions cheap to challenge before they become expensive
code or production incidents.

\section{A practical quality model}
Before choosing tools, ask which of these dimensions are material:

\noindent\begin{tabularx}{\textwidth}{@{}p{0.22\textwidth}X@{}}
\toprule
Dimension & Question to answer \\
\midrule
Correctness & Does the system produce the required result for valid and invalid inputs? \\
Safety and security & What must never happen, and who might deliberately cause it? \\
Reliability & What happens when dependencies, data, or operators behave unexpectedly? \\
Operability & Can a responsible person detect, diagnose, and recover from failure? \\
Changeability & Can a team modify one concern without guessing about unrelated consequences? \\
Human fit & Can people with different abilities, languages, roles, and incentives use and govern it? \\
\bottomrule
\end{tabularx}

These dimensions compete. Stronger durability may cost latency; a highly generic
API may be harder to understand; a fast migration may increase rollback risk.
The engineering response is not to maximize every quality. It is to state the
trade-off and make the consequence visible.

\begin{exercise}
Write five falsifiable properties for a service that accepts payments. For each,
name the boundary at which it is measured and one test or observation that could
show it is false.
\end{exercise}

\section{Chapter summary}
Software engineering is disciplined reasoning under constraints. Begin with
properties rather than technologies, distinguish standards from evidence and
recommendations, draw a deliberate system boundary, and use short feedback
loops to turn assumptions into tests and observations.

\chapter{Requirements, specifications, and stakeholders}

\section{From requests to obligations}
“Make checkout faster” is a request, not yet a requirement. It does not say
which users, which journey, which measurement boundary, or what trade-off is
acceptable. A useful requirement names an actor, an observable behaviour, a
condition, and a threshold or invariant when one matters.

One Ledgerline requirement can be written as:

\begin{quote}
When a caller submits a syntactically valid order with an unused idempotency key,
the service shall persist one accepted order and return its identifier. Repeating
the same request with the same key shall return the original result without
creating another order.
\end{quote}

This statement creates questions: How long is a key retained? What if the first
request timed out after commit? Must the payload match on reuse? What error does a
conflicting payload receive? Requirements become useful when their boundary
cases are discussed rather than hidden.

\section{Stakeholder analysis}
A stakeholder is anyone who can affect, use, constrain, be affected by, or be
held accountable for the system. The person requesting a feature is not always
the person who bears its risks. A small table is often enough to uncover this.

\noindent\begin{tabularx}{\textwidth}{@{}p{0.19\textwidth}p{0.24\textwidth}X@{}}
\toprule
Stakeholder & Interest & Questions \\
\midrule
Customer & A correct, understandable order & What counts as confirmation? Can the order be retried safely? \\
Support & Ability to explain outcomes & Can a person find the request, status, and reason for rejection? \\
Operator & Stable service and recoverable data & What is monitored, backed up, and safe to change? \\
Finance & Reconciled money & Which values are authoritative, and how are corrections audited? \\
Security and privacy & Limited exposure and accountability & Which data is collected, for what purpose, and under whose authority? \\
\bottomrule
\end{tabularx}

Stakeholder analysis does not mean accepting every preference. It makes the
conflict explicit. A customer may want unlimited retention for convenience while
privacy policy requires deletion. An operator may want a safe pause while a
commercial team wants availability. Engineering supplies options and evidence;
governance decides which values prevail.

\section{Functional and quality requirements}
Functional requirements describe behaviour: create an order, cancel it, export a
record. Quality requirements describe conditions on that behaviour: a latency
percentile, a recovery time, an accessibility criterion, or a permitted data
residency. The categories interact. “Returns a result” is incomplete if the
result is too late for the user, cannot be accessed by keyboard, or exposes a
secret in logs.

Use a requirement record with at least:

\begin{enumerate}
  \item identifier and short name;
  \item rationale and stakeholders;
  \item normative statement;
  \item acceptance examples, including invalid cases;
  \item measurable quality target and measurement boundary;
  \item dependencies, assumptions, and unresolved decisions;
  \item owner and review date.
\end{enumerate}

\section{Specifications as executable agreements}
A specification is valuable when independent people can use it to predict the
same behaviour. Examples include an OpenAPI document, a SQL schema, a property
test, a protocol state machine, or a runbook. A prose requirement remains
important for purpose and context, but a machine-checkable representation catches
drift earlier.

Do not make the specification larger than the uncertainty it resolves. A stable
public API deserves explicit compatibility rules. A private function may need
only a type, a test, and a clear name. Excess ceremony creates documents that
are neither read nor updated.

\section{Acceptance and examples}
Examples expose ambiguity. The following set describes the same rule from three
angles:

\begin{itemize}
  \item Given two units at 1,250 cents, the subtotal is 2,500 cents.
  \item Given quantity zero, the request is rejected before persistence.
  \item Given the same idempotency key twice with the same payload, the second
        response refers to the first order.
\end{itemize}

Examples are not a replacement for general rules. They are probes that make the
rule concrete and reveal missing decisions. Property-based tests can then state
the broader invariant: for every valid collection of positive integer prices and
quantities, the total is an integer and equals the sum of each product.

\section{Change and traceability}
Requirements change because the world changes. A useful traceability chain is:
requirement -> design decision -> code -> automated check -> operational signal.
The chain should be lightweight enough to maintain. A commit message or issue
identifier can be sufficient for a small team; regulated work may need a formal
record. The principle is the same: a future reader should be able to ask why a
constraint exists and find evidence.

\begin{tradeoffs}
Detailed upfront specifications reduce some ambiguity but can freeze wrong
assumptions. Pure improvisation learns quickly but can leave invisible contracts.
Use more formality where the cost of misunderstanding is high, and keep the
feedback cycle short where uncertainty is high.
\end{tradeoffs}

\begin{exercise}
Turn “the report page should feel responsive” into a measurable requirement.
State the user journey, the measurement boundary, the percentile and time window,
the test data shape, and one reason the target might need revision.
\end{exercise}

\section{Chapter summary}
Requirements become engineering material when they identify actors, behaviour,
conditions, and evidence. Include quality properties, examine stakeholders who
bear risks, use examples to expose ambiguity, and keep a trace from intent to
operation without creating paperwork that nobody can maintain.

\chapter{Decomposition, abstraction, and boundaries}

\section{The purpose of decomposition}
Decomposition turns a large problem into units that people can understand,
change, and test. It is not merely splitting a file. A good boundary hides a
decision that is likely to change or that should be protected by an invariant.
Parnas's classic discussion treats information hiding as a criterion for module
design, rather than a preference for a particular syntax \cite{parnas1972criteria}.

For Ledgerline, a useful first decomposition is:

\begin{center}
\begin{tikzpicture}[node distance=8mm and 12mm,>=Latex]
  \node[draw,rounded corners,fill=codebg,minimum width=30mm,minimum height=9mm] (api) {HTTP adapter};
  \node[draw,rounded corners,fill=codebg,minimum width=30mm,minimum height=9mm,below=of api] (app) {application service};
  \node[draw,rounded corners,fill=codebg,minimum width=30mm,minimum height=9mm,below=of app] (domain) {order domain};
  \node[draw,rounded corners,fill=codebg,minimum width=30mm,minimum height=9mm,left=of domain] (store) {store port};
  \node[draw,rounded corners,fill=codebg,minimum width=30mm,minimum height=9mm,right=of domain] (events) {event port};
  \node[draw,rounded corners,fill=codebg,minimum width=30mm,minimum height=9mm,below=of store] (sql) {SQL adapter};
  \node[draw,rounded corners,fill=codebg,minimum width=30mm,minimum height=9mm,below=of events] (queue) {queue adapter};
  \draw[->,thick,accent] (api) -- (app);
  \draw[->,thick,accent] (app) -- (domain);
  \draw[->,thick,accent] (domain) -- (store);
  \draw[->,thick,accent] (domain) -- (events);
  \draw[->,thick,accent] (store) -- (sql);
  \draw[->,thick,accent] (events) -- (queue);
\end{tikzpicture}
\end{center}

The arrows represent dependency direction, not necessarily runtime network
traffic. The domain does not need to know whether persistence is SQLite,
PostgreSQL, or an in-memory fake. The adapter owns translation between a
technology's vocabulary and the domain's vocabulary.

\section{Cohesion and coupling}
\term{Cohesion} asks whether the responsibilities inside a unit belong together.
\term{Coupling} asks how much one unit knows about another. High cohesion and
low unnecessary coupling are useful heuristics, not scores to maximize blindly.
Some coupling is valuable: an invariant about money may be kept close to the
database constraint that enforces it. The danger is accidental coupling, such as
a UI importing a database row shape and thereby making a storage migration a UI
change.

A dependency is especially expensive when it crosses:

\begin{itemize}
  \item ownership boundaries, because coordination becomes necessary;
  \item time boundaries, because queues and retries introduce duplication or delay;
  \item trust boundaries, because input must be authenticated and validated;
  \item data-lifetime boundaries, because privacy and deletion rules differ;
  \item release boundaries, because two versions must coexist.
\end{itemize}

\section{Abstraction and the cost of being wrong}
An abstraction removes detail so that a person can reason about a larger unit.
It also hides information. An abstraction is successful when the hidden detail
can vary without surprising its clients. It is unsuccessful when the abstraction
promises a simple concept but leaks the complexity through exceptions, latency,
ordering, or resource usage.

The phrase “repository” can hide a database, but it cannot make a database
transaction disappear. If callers assume that `save` is atomic while the
implementation performs three independent writes, the abstraction lies. Make
important semantics explicit: \texttt{save\_order\_and\_outbox\_event} communicates more
than `save` if atomicity is essential.

\section{Designing a boundary}
For each proposed module, write four sentences:

\begin{enumerate}
  \item Its responsibility is ...
  \item It guarantees ...
  \item It may change without requiring ...
  \item It must not know ...
\end{enumerate}

Then test the boundary against a plausible change. If changing the tax rate
requires editing HTTP parsing, SQL queries, and every unit test, the domain rule
is scattered. If a tax module cannot be tested without a network and a database,
the boundary is probably too entangled.

\section{When a boundary is too early}
Premature decomposition creates indirection before the shape of the problem is
understood. Five interfaces around five lines of code can make a small system
harder to read. Start with seams that protect a real decision: an external API,
a security check, a durable write, or a volatile policy. Refactor toward clearer
boundaries as evidence accumulates. The repository's Python example keeps the
domain function small because the purpose is visible; it does not manufacture a
framework around it.

\begin{exercise}
Take a single function that validates and persists an order. List three changes
that could occur independently. Design boundaries that isolate those changes,
then name one cost introduced by each boundary.
\end{exercise}

\section{Chapter summary}
Decompose around changing decisions and protected invariants. Use information
hiding, cohesion, and coupling as reasoning tools. Treat abstractions as
contracts about semantics, not just shorter names, and delay indirection when it
does not yet protect a real change.

\chapter{APIs, data modelling, and persistence}

\section{An interface is a promise over time}
An \term{API} is more than a URL or method signature. It defines inputs, outputs,
errors, identity, authorization, ordering, consistency, and compatibility. A
client can depend on behaviour that was never written down, so an API should
make its important promises deliberate.

For an HTTP endpoint, document at least:

\begin{itemize}
  \item the resource and method semantics;
  \item accepted representations and validation rules;
  \item success and error status codes with stable error shapes;
  \item authentication, authorization, rate limits, and sensitive fields;
  \item retry and idempotency behaviour;
  \item pagination, ordering, time zones, and version-compatibility rules.
\end{itemize}

HTTP is standardized as a stateless application-level protocol with shared
semantics across versions; the standard does not dictate your business resource
model \cite{ietf2022http}. This distinction matters. A `POST /orders` contract
is yours, while the semantics of status codes and method properties are governed
by the HTTP specification.

\section{Identifiers and time}
Identifiers should communicate less than they need to. An opaque identifier
reduces accidental coupling to storage order and makes enumeration harder, but
it does not replace authorization. If a client needs chronological ordering,
provide an explicit timestamp or cursor rather than asking it to decode an ID.

Time has at least three meanings: an instant in a global timeline, a local
calendar representation, and a duration measured by a monotonic clock. Store
instants in an unambiguous form, carry time-zone information when presenting
calendar values, and use a monotonic clock for elapsed-time measurements. Never
use wall-clock time to decide whether a timeout duration has elapsed.

\section{Relational modelling}
A relational schema is an executable set of invariants. In the example schema,
positive quantities are checked in the database even though application code
also validates them. Defence in depth is useful because imports, migrations,
administrative scripts, and future services can bypass one validation path.

The order total is stored as integer cents. Decimal money is a domain choice;
the important point is that binary floating-point approximations are not silently
used for a value whose legal meaning is exact. The database view recalculates the
sum of line items. In a production design, decide whether the stored total is a
snapshot, a cache, or an authoritative ledger value, and reconcile it explicitly.

\begin{lstlisting}[language=SQL,caption={A constraint is part of the data contract.}]
CREATE TABLE order_items (
    order_id TEXT NOT NULL REFERENCES orders(id),
    sku TEXT NOT NULL CHECK (length(trim(sku)) > 0),
    quantity INTEGER NOT NULL CHECK (quantity > 0),
    unit_price_cents INTEGER NOT NULL CHECK (unit_price_cents > 0),
    PRIMARY KEY (order_id, sku)
);
\end{lstlisting}

\section{Transactions and consistency}
A transaction groups changes so that the database presents an allowed state at
its commit boundary. The exact isolation guarantees vary by database engine and
configuration. Do not infer serializable behaviour from the word “transaction.”
Ask what anomalies are possible: dirty reads, lost updates, write skew, or
phantoms. Then choose a constraint, lock, isolation level, or retry strategy that
matches the invariant.

The outbox pattern is useful when a durable state change must cause a message.
Write the domain change and an outbox row in the same local transaction. A
separate publisher reads unclaimed outbox rows and sends them with at-least-once
delivery. Consumers must therefore be idempotent. The pattern does not create
exactly-once effects across arbitrary systems; it creates a recoverable handoff.

\section{Schema evolution}
A safe migration is compatible with the versions that may coexist during rollout.
The expand-and-contract sequence is a common strategy:

\begin{enumerate}
  \item add a nullable or additive representation;
  \item deploy code that can read both and writes the new form;
  \item backfill and validate existing data;
  \item switch readers and enforce the new invariant;
  \item remove the old representation only after rollback is no longer needed.
\end{enumerate}

This costs temporary complexity. The alternative, a single destructive change,
may be reasonable for a private database with a verified restore path and a
maintenance window. The decision depends on data value, downtime tolerance,
rollout topology, and rollback feasibility.

\begin{exercise}
Design an endpoint for retrieving a customer's orders. Specify pagination,
ordering, authorization, deletion behaviour, and what happens when a client
repeats a request after a timeout.
\end{exercise}

\section{Chapter summary}
APIs are long-lived behavioural contracts. Model identity, time, errors, retries,
and compatibility explicitly. Use database constraints as executable invariants,
understand the actual transaction guarantees, and evolve schemas in a way that
respects the versions present during deployment.


\chapter{Concurrency, networking, and distributed systems}

\section{Concurrency is about interleaving}
Concurrency means multiple activities can make progress during overlapping
intervals. Parallelism is simultaneous execution. A single-threaded event loop
can be concurrent without being parallel; multiple CPU cores can be parallel.
The distinction matters for reasoning, resource limits, and performance claims.

The basic danger is an interleaving that violates an invariant. Suppose two
requests both read an account balance of 100 and both decide that a withdrawal
of 80 is allowed. If each writes 20, the final balance looks plausible but the
system has accepted 160 of withdrawals. A transaction, an atomic conditional
update, or a serialized command can protect the rule.

\begin{keyidea}
Do not ask only whether each operation is correct in isolation. Ask which
interleavings are possible, which state is shared, and which invariant must hold
after every visible transition.
\end{keyidea}

\section{Memory and message boundaries}
In one process, threads share memory and require synchronization. Across
processes, machines, or organizations, messages cross a boundary and can be
delayed, duplicated, reordered, or lost. A network timeout does not tell the
caller whether the remote operation failed or succeeded but its response was
lost. This is the source of many duplicate writes.

A robust request path defines:

\begin{itemize}
  \item a deadline propagated through every downstream call;
  \item a bounded retry policy only for operations that are safe to repeat;
  \item an idempotency key or deduplication record for retried effects;
  \item a clear failure result when the outcome is unknown;
  \item telemetry that distinguishes timeout, refusal, overload, and validation.
\end{itemize}

Retries multiply load during failure. Exponential backoff with jitter spreads
attempts, but it is not a cure for an overloaded dependency. A caller should
retry only when the error is plausibly transient and the total deadline still
allows useful work. A circuit breaker can stop calls to a failing dependency,
but it introduces its own state machine and false-open trade-off.

\section{Ordering and causality}
Wall-clock timestamps are not a reliable global ordering in a distributed system.
Clocks differ and can move. Lamport's work on logical clocks shows how a system
can represent a partial order of events without pretending that every event has
one exact global time \cite{lamport1978time}. In practice, choose the smallest
ordering guarantee that the product needs: per-account sequence numbers,
partition-local order, or an explicit version check.

\section{Consistency is a product decision}
An eventually consistent view can be correct for a feed whose purpose is
discovery, but unacceptable for a bank balance or an authorization decision.
“Strong” and “eventual” are not complete specifications. State what a reader may
observe after a write, which conflicts can occur, and how a user can repair or
understand them.

The CAP theorem is frequently reduced to a slogan about choosing two of three
properties. The more useful interpretation is narrower: under a network
partition, a distributed data system cannot simultaneously provide a particular
strong consistency guarantee and uninterrupted availability for every operation.
The actual design space includes latency, scope of consistency, failure model,
and recovery strategy \cite{brewer2012cap}.

\section{Queues and workers}
Queues decouple producers and consumers in time. They absorb bursts and allow a
worker pool to process work at a controlled rate. They also create backlog,
duplicate delivery, poison messages, and delayed user-visible state. A queue is
not a free buffer; it is a new stateful subsystem.

For a message handler, specify:

\begin{enumerate}
  \item what makes a message valid;
  \item the acknowledgment point;
  \item the effect of a crash before and after acknowledgment;
  \item maximum attempts and dead-letter handling;
  \item deduplication identity and retention;
  \item how operators replay, quarantine, or delete a message safely.
\end{enumerate}

\section{A small state machine}
Model a payment or order lifecycle as a state machine before implementing it:

\begin{center}
\begin{tikzpicture}[node distance=20mm,>=Latex]
  \node[draw,rounded corners,fill=codebg,minimum width=25mm,minimum height=9mm] (new) {new};
  \node[draw,rounded corners,fill=codebg,minimum width=25mm,minimum height=9mm,right=of new] (accepted) {accepted};
  \node[draw,rounded corners,fill=codebg,minimum width=25mm,minimum height=9mm,right=of accepted] (paid) {paid};
  \node[draw,rounded corners,fill=codebg,minimum width=25mm,minimum height=9mm,below=of accepted] (cancelled) {cancelled};
  \draw[->,thick,accent] (new) -- node[above,font=\scriptsize]{validate} (accepted);
  \draw[->,thick,accent] (accepted) -- node[above,font=\scriptsize]{capture} (paid);
  \draw[->,thick,accent] (accepted) -- node[right,font=\scriptsize]{cancel} (cancelled);
\end{tikzpicture}
\end{center}

The diagram makes illegal transitions visible. It also raises operational
questions: what if capture succeeds but the response times out? The answer may
be a reconciliation state rather than an optimistic retry.

\begin{exercise}
For a queue consumer that sends an email, list the states before and after a
crash at each of these points: before send, after send, before acknowledgment,
and after acknowledgment. Choose a deduplication strategy and state its limits.
\end{exercise}

\section{Chapter summary}
Concurrency is reasoning about interleavings; distribution adds uncertain
messages and partial failure. Define deadlines, retryability, idempotency,
ordering, consistency, and recovery. Treat queues and state machines as explicit
parts of the product rather than invisible implementation details.


\chapter{Architecture styles and trade-offs}

\section{Architecture is a set of consequential decisions}
Architecture describes the significant structures, boundaries, and constraints
that shape a system. A diagram is useful only when it helps someone decide,
implement, operate, or change the system. A box labelled “service” is not an
architecture unless its responsibility, communication, data ownership, and
failure behaviour are understood.

\section{Common styles}
\noindent\begin{tabularx}{\textwidth}{@{}p{0.23\textwidth}p{0.26\textwidth}X@{}}
\toprule
Style & Helps when & Costs and failure modes \\
\midrule
Modular monolith & One team needs strong local transactions and simple deployment & A shared process can become tangled; scaling and isolation are coarse. \\
Layered system & Responsibilities have a stable vertical flow & Layers become pass-through plumbing or leak abstractions. \\
Ports and adapters & Domain logic should remain independent of infrastructure & Too many interfaces obscure a small system; mappings must be maintained. \\
Event-driven & Producers and consumers need temporal decoupling & Ordering, retries, schemas, and debugging become distributed concerns. \\
Microservices & Teams need independent ownership and deployment at meaningful boundaries & Network failure, data ownership, versioning, and operational cost multiply. \\
\bottomrule
\end{tabularx}

None is the universally mature choice. A microservice is not automatically more
scalable, and a monolith is not automatically a prototype. The decision should
be driven by boundaries of change, ownership, failure isolation, deployment
independence, data consistency, and capacity rather than by fashion.

\section{Architecture decision records}
An architecture decision record (ADR) keeps a decision close to its reasons:

\begin{enumerate}
  \item Context: what forces and constraints exist?
  \item Decision: what will we do?
  \item Alternatives: what serious options were considered?
  \item Consequences: what becomes easier, harder, or newly risky?
  \item Status: proposed, accepted, superseded, or rejected.
\end{enumerate}

An ADR is not a vote and not a substitute for measurement. It is a memory aid.
If a decision was made under an assumption about traffic, team size, or a vendor,
record the assumption and a trigger for revisiting it.

\section{Data ownership}
In a distributed architecture, the hardest boundary is often data. If two
services write the same tables, ownership is nominal. If they have separate
stores, cross-service transactions become workflows with compensation or
reconciliation. The right choice depends on consistency requirements and
failure tolerance.

For Ledgerline, keep order state in one database while the domain is small. A
separate payment service may make sense when legal responsibility, risk controls,
or independent scaling require it. Splitting it merely to create more boxes
would make failures harder to recover from without solving a real constraint.

\section{Architecture fitness functions}
A \term{fitness function} is a repeatable check for an architectural property.
Examples:

\begin{itemize}
  \item a dependency rule that prevents the domain package importing HTTP code;
  \item a migration test that proves an old and new application version can coexist;
  \item a load test that measures p95 latency at a stated arrival rate;
  \item a security test that rejects access across a tenant boundary;
  \item an accessibility test that checks keyboard navigation for a critical flow.
\end{itemize}

Fitness functions do not prove the architecture is good. They catch regression
against properties the team has decided matter. A metric without a decision
attached can become decorative or actively misleading.

\section{Architecture review}
Review the smallest set of questions that could change the design:

\begin{itemize}
  \item What are the critical user journeys and failure consequences?
  \item Which state is authoritative and which copies are disposable?
  \item What happens when every dependency is slow, unavailable, or inconsistent?
  \item How is a change released and rolled back while old clients remain?
  \item Which trust boundaries, secrets, personal data, and audit obligations exist?
\end{itemize}

\begin{tradeoffs}
Centralization simplifies coordination and can create a bottleneck. Distribution
can improve isolation and team autonomy and can also multiply failure modes. The
correct architecture is the smallest one that meets the real constraints with
known consequences.
\end{tradeoffs}

\begin{exercise}
Choose a modular monolith, a service split, or an event-driven design for an
application that has one team, 20 requests per second, a strict audit trail, and
one database. Defend the choice, then identify the first signal that would make
you reconsider it.
\end{exercise}

\section{Chapter summary}
Architecture is about boundaries and consequences. Select a style for ownership,
failure, change, data, and capacity reasons. Record assumptions in ADRs and turn
important architectural properties into automated or observable fitness checks.

\chapter{Object-oriented, functional, procedural, and data-oriented design}

\section{Paradigms are lenses}
A programming paradigm emphasizes certain ways of representing state,
behaviour, and composition. Object-oriented design groups data with operations
and uses encapsulation and polymorphism. Functional design emphasizes expressions,
immutable values, and explicit transformation. Procedural design emphasizes
ordered commands and control flow. Data-oriented design arranges data for the
access patterns of the hardware and workload.

Real systems combine them. A functional core can sit inside an object-oriented
application; a procedural migration script can protect a data-oriented storage
layout. The question is whether the chosen representation makes the important
invariants and costs visible.

\section{Encapsulation and state}
Encapsulation means clients depend on a stable set of operations rather than
directly manipulating representation. It helps when the invariant is local and
the state has a meaningful lifecycle. It becomes harmful when a class is merely a
namespace for unrelated functions or when inheritance creates hidden coupling.

Prefer composition when behaviour varies independently. Inheritance can express
a genuine substitutability relation, but sharing code is not by itself a reason
to create a subtype. A payment gateway that supports `authorize` and `capture`
may be a port with several adapters; it need not form a deep class hierarchy.

\section{Functional core, imperative shell}
Pure functions are easier to test because their result depends on explicit input.
Side effects are still necessary for useful software. A practical split is:

\begin{itemize}
  \item imperative shell: read HTTP, clock, randomness, database, and network;
  \item functional core: validate, calculate totals, decide transitions, and
        construct commands from explicit values.
\end{itemize}

The Python \texttt{total\_cents} example follows this approach. It accepts data, applies
rules, and has no network or database dependency. The shell decides when and how
to call it. This does not eliminate integration tests; it makes the rule-level
tests fast and diagnostic.

\section{Procedures and data flow}
Procedural code can be the clearest form for a workflow. A migration or a shell
script is often easier to audit as a sequence of explicit steps than as a graph
of objects. The risk is accidental shared state and unclear error handling. Keep
preconditions, postconditions, and cleanup paths visible.

\section{Data-oriented thinking}
Performance depends on the whole path: algorithm, data layout, allocation,
cache behaviour, I/O, concurrency, and workload shape. Data-oriented design
starts with the hot operations and organizes data so they do not repeatedly
load irrelevant fields. It is useful in simulations, games, analytics, and
high-throughput infrastructure; it can make ordinary business code less
readable when applied without a measured bottleneck.

\section{Choosing a representation}
Use these questions:

\begin{enumerate}
  \item What state must be protected by an invariant?
  \item Which effects are unavoidable, and can they be isolated?
  \item Does behaviour vary by a stable capability or by a changing data value?
  \item Is the workload limited by CPU, memory, allocation, I/O, or coordination?
  \item What will a future reader need to change safely?
\end{enumerate}

\noindent\begin{tabularx}{\textwidth}{@{}p{0.24\textwidth}X@{}}
\toprule
Situation & Often useful starting point \\
\midrule
Local lifecycle invariant & Encapsulation around a small state machine \\
Many transformations and rules & Pure functions and immutable values \\
Explicit workflow or migration & Procedural sequence with checked steps \\
Measured hot path and predictable data & Data-oriented layout and profiling \\
\bottomrule
\end{tabularx}

\begin{exercise}
Represent an order cancellation in two ways: a mutable object and a pure
transition function. For each, state where invalid state is prevented and where
side effects occur. Which representation gives the clearer test failure?
\end{exercise}

\section{Chapter summary}
Paradigms are tools for making state, effects, and cost visible. Combine them
deliberately. Prefer the simplest representation that protects the important
invariants and fits the measured workload, and avoid turning stylistic taste into
an architectural law.

\chapter{Version control, review, and collaborative development}

\section{Version control as a reasoning tool}
Version control records a sequence of states and the intent behind changes. Git
is a distributed revision control system; its commands let a team compare,
select, combine, and recover work. The commands are not the process. A process
is good when it makes change reviewable, integrates often enough to expose
conflicts, and preserves a trustworthy path from request to released state.

A useful commit is coherent: one idea, a message that explains why, and tests
that support the changed behaviour. Small does not mean trivial. A database
migration and the code that requires it may belong in one change if separating
them would create an invalid intermediate state.

\section{Code review}
Review is a risk-reduction activity, not a search for stylistic agreement. Read
the requirement, then the design, then the diff and its tests. Ask:

\begin{itemize}
  \item Is the behaviour correct for valid, invalid, and repeated inputs?
  \item Does the change preserve data, security, compatibility, and performance
        properties that matter?
  \item Are failures observable and recoverable?
  \item Is the new abstraction justified by a real boundary?
  \item Can a future maintainer understand the reason and remove obsolete code?
\end{itemize}

The Git project's own review guidance emphasizes correctness, clarity, and the
importance of reviewing test coverage, while also noting that its advice is not
binding for every project \cite{git2026review}. This is a good example of a
professional recommendation rather than a law.

\section{Review comments}
Prefer comments that identify risk and invite a decision:

\begin{quote}
The retry happens after the database write. If the response is lost, the caller
can submit again and create a second order. Can we store the idempotency key with
the result, or document why duplicate creation is acceptable here?
\end{quote}

This is better than “please refactor.” It connects a concrete path to a property.
Distinguish blocking correctness or security issues from optional suggestions.
Resolve disagreements by testing assumptions or asking the owner of the
requirement to decide, not by escalating the loudest preference.

\section{Documentation that earns its keep}
Documentation serves different readers:

\noindent\begin{tabularx}{\textwidth}{@{}p{0.24\textwidth}X@{}}
\toprule
Document & Primary question \\
\midrule
README & How do I run, test, and orient myself? \\
API reference & What can I call, and what happens? \\
ADR & Why is the architecture this way? \\
Runbook & What should an operator do in a known situation? \\
Comment & Why is this local code non-obvious? \\
\bottomrule
\end{tabularx}

Keep a document near the code it describes, automate examples where possible,
and delete claims that no longer hold. A short, verified runbook is more useful
than a comprehensive page of guesses.

\section{Branching and integration}
Long-lived branches defer integration risk. Short-lived branches and continuous
integration expose conflicts while the author still remembers the change. This
does not require one branching model. A team may use trunk-based development,
release branches, or a hybrid when release governance demands it. What matters is
that the model matches deployment cadence and that it is possible to test the
state being reviewed.

\section{Psychological safety and accountability}
Review quality depends on the social system. People need permission to report a
risk, question a requirement, and admit uncertainty. Psychological safety does
not mean low standards; it makes high standards discussable. Accountability is
stronger when the system records decisions and evidence instead of assigning
blame for every defect.

\begin{exercise}
Write three review comments on a hypothetical change that adds a cache. One
should concern correctness, one privacy or security, and one operability. Mark
which are blocking and why.
\end{exercise}

\section{Chapter summary}
Use version control to preserve intent and reviewable history. Review behaviour,
risk, tests, and operational consequences rather than personal style. Keep
documentation purpose-specific, integrate frequently enough to expose risk, and
build a social environment where uncertainty can be reported early.

\chapter{Testing, analysis, debugging, and proof}

\section{Testing as evidence}
A test is an experiment with a programmed oracle. It does not prove that all
possible behaviour is correct. It provides evidence about a stated input space,
environment, and observation. The strongest test suite is not the one with the
largest line count; it is the one that catches important failures and explains
them quickly.

\noindent\begin{tabularx}{\textwidth}{@{}p{0.22\textwidth}p{0.25\textwidth}X@{}}
\toprule
Level & Best at finding & Main cost \\
Unit & Local rule and edge-case defects & Can miss wiring, configuration, and real dependencies \\
Component & Contract and integration defects & More setup and slower diagnostics \\
End to end & Broken user journeys and deployment assumptions & Slow, brittle, and often hard to localize \\
Property or model & Broad invariant violations and state-machine bugs & Requires a good model and failure shrinking \\
\bottomrule
\end{tabularx}

The levels are not a pyramid law. A small system with a strong component test may
need few end-to-end tests. A safety-critical subsystem may need formal analysis.
Choose the mix from failure cost, change rate, and observability.

\section{Test design}
For each behaviour, identify:

\begin{enumerate}
  \item a representative valid input;
  \item boundary values and malformed input;
  \item a state transition or dependency failure;
  \item the observable result and any invariant that must remain true.
\end{enumerate}

Avoid asserting incidental representation. If a client only needs an order ID
and status, do not make it depend on database column order or an internal class.
Tests are also consumers of an interface; over-specified tests make refactoring
expensive without increasing confidence.

\section{Static analysis and types}
Static analysis checks source without executing every path. A type checker can
catch incompatible values; a linter can catch suspicious constructs; a dataflow
tool can identify possible taint or resource leaks. False positives impose cost,
so tune rules and provide an explicit escape hatch with a reason. A passing
static check is evidence, not a proof of correctness.

Types are especially useful when they encode meaningful distinctions. A string
for a customer ID and a string for a password are technically interchangeable but
not semantically interchangeable. A wrapper type or validation boundary can
reduce accidental mixing. Do not create types that obscure a simple domain rule.

\section{Debugging as hypothesis testing}
Debugging is a search through competing explanations. Start with a reproducible
failure and preserve the original evidence. Write the expected and actual
behaviour, then reduce the input or environment while keeping the failure.

A practical loop is:

\begin{enumerate}
  \item observe: capture logs, traces, inputs, versions, and timing;
  \item hypothesize: name a mechanism that could explain the observation;
  \item predict: identify what else should be true if the hypothesis is correct;
  \item test: run the smallest discriminating experiment;
  \item change: fix the cause or update the model, then add a regression check.
\end{enumerate}

Changing several things at once destroys information. The exception is an active
incident where stabilization is more important than diagnosis; record temporary
changes so they can be revisited.

\section{Formal methods in proportion}
Formal methods use mathematical models to specify or analyze systems. They range
from assertions and contracts to model checking, theorem proving, and formally
verified compilation. The benefit is strongest when a small, critical state space
can be modeled precisely. The cost is building and maintaining the model and
connecting it to implementation.

An invariant can be simple:

\[
  \forall i \in \text{order items}: i.quantity > 0 \land i.unit\_price\_cents > 0.
\]

That statement can be enforced by application validation, a database constraint,
and tests. More advanced tools can explore concurrent transitions. Use formal
analysis where it reduces a material risk, not as decoration.

\section{Mutation and negative evidence}
Mutation testing changes code in small ways and asks whether the tests fail. A
surviving mutation points to a missing assertion, but mutation scores are not a
universal quality metric. Similarly, a security test that finds no vulnerability
does not establish security; it establishes only that the tested attack and
observation did not reveal one.

\begin{exercise}
Create a test matrix for an order endpoint with dimensions of authentication,
payload validity, duplicate key, database failure, and client timeout. Choose the
smallest set of tests that covers each important interaction and explain what is
left untested.
\end{exercise}

\section{Chapter summary}
Tests provide scoped evidence. Combine unit, component, journey, property, and
static checks according to risk. Debug by preserving evidence and testing
hypotheses. Use types, assertions, and formal models when they clarify and reduce
important risk rather than because a technique is fashionable.

\chapter{Reliability, resilience, scalability, and capacity}

\section{Define reliable for whom}
Reliability is not “the server did not crash.” It is the probability that a
system performs its required function for a stated user, time, and environment.
An API can be available while returning stale or incorrect data; a batch job can
be correct while missing its deadline. Define the service indicator before
choosing a target.

Useful indicators include successful request ratio, freshness, correctness rate,
queue age, and latency percentiles. The percentile must include a population,
window, and boundary. “p95 under 200 ms” is incomplete unless it says which
requests, which route, whether retries are included, and where time starts and
ends.

\section{SLOs and error budgets}
A service-level objective (SLO) is a target such as “99.9 percent of eligible
requests in a rolling 30-day window succeed within the defined latency bound.”
The remaining error budget is the permitted failure fraction. It is a decision
tool: when the budget is healthy, a team may accept a risky improvement; when it
is exhausted, stability work has priority. Google describes SLOs as a way to
state and measure desired reliability \cite{google2016sre}.

The arithmetic is simple. For a 30-day window, 99.9 percent availability allows
about 43.2 minutes of unavailable time if the measurement is continuous. The
number is not a promise: maintenance, exclusions, sampling, and user impact
change the interpretation. State the measurement method.

\section{Failure modes and resilience}
Resilience is the ability to continue useful operation or recover after failure.
List failure modes by layer:

\begin{itemize}
  \item input: malformed, oversized, duplicated, or adversarial;
  \item process: crash, memory exhaustion, deadlock, resource leak;
  \item dependency: timeout, partial response, incompatible version;
  \item data: corruption, lag, missing migration, unrecoverable deletion;
  \item human and organizational: bad release, unclear ownership, unavailable expert.
\end{itemize}

For each, choose a response: prevent, detect, degrade, retry, isolate, fail
closed, fail open, or recover. “Fail gracefully” is not a response until the
degraded behaviour is named.

\section{Scalability and capacity}
Scalability is how resource use and performance change as workload grows. It is
not synonymous with speed. Capacity planning starts with a workload model:
arrival rate, request mix, payload size, concurrency, data growth, and burst
shape. Measure service time and a bottleneck resource. Little's Law gives a
useful relationship for a stable queueing system:

\[
  L = \lambda W,
\]

where $L$ is average work in the system, $\lambda$ is average arrival rate, and
$W$ is average time in the system. The relationship does not by itself predict
tail latency or behaviour during overload, but it exposes inconsistent capacity
claims.

\section{Load, stress, and performance tests}
A performance test is only meaningful with a stated workload, environment,
warm-up, data shape, concurrency, and measurement. A benchmark that reports
average latency may hide a failing tail. A stress test beyond capacity can reveal
recovery behaviour, but it should not be run against a production dependency
without explicit authorization and safeguards.

Optimization should follow a measured bottleneck. Cache design changes
correctness and invalidation; parallelism changes ordering and contention; a
database index changes write cost and storage. Record the baseline, the change,
and the new workload result.

\section{Backups and recovery}
A backup that has never been restored is an untested hope. Define recovery point
objective (RPO) and recovery time objective (RTO), encrypt backups, control
access, test restoration, and document dependencies such as keys and schema
versions. A replication system can copy corruption or accidental deletion, so it
is not a substitute for historical recovery points.

\begin{exercise}
Write an SLO for an order-creation endpoint, calculate its monthly error budget,
and name three events that should consume it. Then specify an RPO and RTO for
orders and justify them in terms of user and business impact.
\end{exercise}

\section{Chapter summary}
Reliability starts with a defined user-visible indicator. Use SLOs to make trade-
offs explicit, model failure modes, plan capacity from workload, test performance
with reproducible conditions, and treat restoration as a capability that must be
exercised.


\chapter{Security engineering, privacy, and threat modelling}

\section{Security is a system property}
Security engineering protects assets against unwanted actions under a stated
threat model. It includes confidentiality, integrity, availability, authenticity,
accountability, and privacy. A library patch cannot compensate for an endpoint
that authorizes the wrong tenant or logs credentials.

Start by naming assets and consequences:

\noindent\begin{tabularx}{\textwidth}{@{}p{0.23\textwidth}p{0.26\textwidth}X@{}}
\toprule
Asset & Adversary or accident & Consequence \\
\midrule
Order data & Cross-tenant read or export error & Confidentiality and legal exposure \\
Payment state & Forged callback or replay & Financial loss and reconciliation failure \\
Credentials & Log leak or phishing & Account takeover \\
Deployment pipeline & Malicious dependency or stolen token & System-wide compromise \\
\bottomrule
\end{tabularx}

\section{Threat modelling}
A lightweight threat model records actors, assets, entry points, trust boundaries,
and abuse cases. STRIDE is one categorization technique; attack trees, misuse
cases, and data-flow diagrams are alternatives. The technique is useful when it
changes a design or test. Listing threats without owners or mitigations is a
worksheet, not security engineering.

For an order API, ask:

\begin{enumerate}
  \item Can an unauthenticated caller reach it?
  \item How is tenant ownership derived and checked on every read and write?
  \item Can a caller replay a payment callback or idempotency key?
  \item Are quantities, IDs, URLs, files, and error messages bounded and encoded?
  \item Which secrets cross the process, logs, build system, and backups?
\end{enumerate}

\section{Secure defaults and input handling}
Default-deny authorization, least privilege, parameterized database queries,
context-appropriate output encoding, safe error messages, bounded input, and
secure transport are durable principles. “Validate input” is incomplete: validate
syntax at the boundary, enforce domain rules in the domain, and enforce critical
invariants in storage. Canonicalize before comparing where multiple encodings
could represent the same resource.

Authentication identifies a principal; authorization decides whether that
principal may perform a specific action on a specific resource. An authenticated
user is not automatically authorized to read every order. Check authorization
server-side, close to the resource access, and test both allowed and denied paths.

The OWASP ASVS provides a versioned basis for testing web application security
controls and requirements; its identifiers can change between versions, so cite
the exact version when using a requirement \cite{owasp2025asvs}.

\section{Secrets and cryptography}
Do not store secrets in source control, client bundles, ordinary logs, or error
messages. Use a secret manager or protected environment injection, rotate keys,
limit access, and record a plan for compromise. Choose well-reviewed cryptographic
primitives through a maintained library. Do not invent encryption or password
hashing schemes. Encryption does not solve authorization, key management, or
metadata leakage.

\section{Privacy by design}
Privacy includes purpose limitation, data minimization, access control,
retention, deletion, transparency, and accountability. A personal-data field
should have a reason, an owner, a retention rule, and an access path. Test exports
and deletion across caches, replicas, logs, backups, analytics, and derived data.
Legal obligations vary by jurisdiction; this chapter offers engineering
controls, not legal advice.

\section{Supply-chain security}
Record dependency versions and integrity, review transitive risk, protect build
credentials, produce provenance where possible, and make the release artifact
reproducible enough to investigate. NIST's SSDF provides a high-level vocabulary
for integrating secure development practices into a lifecycle \cite{nist2022ssdf}.
As of this edition's research date, NIST had also published an initial public
draft for a revision; a draft should not be represented as a final standard
\cite{nist2025ssdfdraft}.

\section{Security verification}
Use layered evidence: code review, static checks, dependency review, unit and
integration tests, dynamic testing, configuration checks, secret scanning, and
authorized penetration testing. Keep a vulnerability response path with
severity, owner, affected versions, remediation, disclosure, and verification.
No checklist proves absence of vulnerabilities. The aim is to reduce likely
failure, limit impact, and improve recovery.

\begin{exercise}
Draw a data-flow diagram for a multi-tenant order API. Mark trust boundaries and
write five abuse cases. For each, choose one preventive control and one test that
would detect a regression.
\end{exercise}

\section{Chapter summary}
Security is a property of architecture, code, data, people, and operations.
Model assets and adversaries, separate authentication from authorization, use
secure defaults and layered validation, minimize and govern personal data, and
keep standards versioned and evidence-based.

\chapter{CI/CD, containers, cloud systems, and infrastructure as code}

\section{Delivery is a controlled experiment}
Continuous integration (CI) makes changes build and test together. Continuous
delivery keeps a releasable artifact available; continuous deployment may
automatically release it after required checks. These are related practices, not
synonyms. The useful question is how quickly and safely the team can move from a
reviewed change to an observable, reversible production state.

A pipeline should answer:

\begin{enumerate}
  \item What exact source revision produced this artifact?
  \item Which tests, analyses, and security checks ran?
  \item Which configuration is injected at deployment rather than baked into the artifact?
  \item How is the artifact promoted, observed, rolled back, or rolled forward?
  \item Who or what is authorized to perform each transition?
\end{enumerate}

\section{Build once, configure at runtime}
Build a versioned artifact once and promote that artifact through environments.
Rebuilding separately for staging and production can produce different bytes and
invalidate test evidence. Configuration such as endpoints and feature flags may
vary; code and dependencies should not silently vary.

Pin dependencies and base images to a reviewed policy. “Latest” is convenient
for a tutorial and dangerous for reproducibility. Updates still need to happen;
the point is to make them explicit and observable.

\section{Containers}
A container packages a process and its user-space dependencies while sharing the
host kernel. It is not a complete security boundary by itself. Use a minimal
image, run as a non-root user, drop unnecessary capabilities, make the filesystem
read-only where feasible, bound resources, and avoid embedding secrets. The
example Compose file demonstrates several of these controls.

Docker's current documentation describes the Compose Specification as the
recommended format for defining services, networks, and volumes; the format is
implemented by modern Compose CLI versions \cite{docker2026compose}. Commands
and implementation details remain version-dependent, so CI should verify the
toolchain it expects.

\section{Infrastructure as code}
Infrastructure as code makes desired infrastructure reviewable and repeatable.
It does not make infrastructure deterministic automatically. Providers have
eventual consistency, external state, secrets, quotas, and destructive changes.

Treat infrastructure changes like application changes:

\begin{itemize}
  \item plan and review the diff;
  \item separate creation, migration, and deletion risk;
  \item protect state and limit its sensitive contents;
  \item test modules in a disposable environment;
  \item document manual recovery when automation is unavailable.
\end{itemize}

\section{Deployment strategies}
\noindent\begin{tabularx}{\textwidth}{@{}p{0.22\textwidth}p{0.27\textwidth}X@{}}
\toprule
Strategy & Benefit & Risk or requirement \\
\midrule
Rolling & Gradual replacement with one pool & Old and new versions must coexist \\
Blue-green & Fast switch between two environments & Doubles capacity and complicates state \\
Canary & Evidence from a small traffic slice & Requires good segmentation and rollback signals \\
Feature flag & Decouples code release from user exposure & Flag lifecycle and combinatorial states \\
\bottomrule
\end{tabularx}

Rollback is not always safe. A database migration, irreversible message, or
external side effect can make old code unable to run. Design forward-compatible
changes and a compensating action where possible.

\section{CI workflow example}
The repository workflow uses pinned major action versions, a fresh checkout,
separate example and PDF jobs, and an uploaded artifact. In a real project, add
permissions with least privilege, dependency review, test result retention,
protected environments, and secret handling appropriate to the platform.

\begin{exercise}
Design a release pipeline for a service with a schema migration and a queue
consumer. Mark which steps can be automatic, which require approval, and what
signal triggers rollback or a pause.
\end{exercise}

\section{Chapter summary}
Delivery should produce traceable, reproducible, observable transitions. Build
once, configure deliberately, harden containers, review infrastructure changes,
and choose deployment strategies based on compatibility and rollback reality.

\chapter{Observability, operations, incidents, and postmortems}

\section{From running to knowing}
A running process is not necessarily a useful service. Operations asks whether
the service is delivering its intended outcome and what a responsible person can
do when it is not. Observability is the ability to infer internal state from
externally emitted signals. OpenTelemetry defines a vendor-neutral set of
telemetry concepts and uses normative terms such as MUST and SHOULD inside its
specification \cite{opentelemetry2026spec}.

\section{Signals}
\term{Metrics} summarize values over time: request count, error count, latency,
queue age, saturation. \term{Logs} record discrete events with context.
\term{Traces} connect work across process and network boundaries. Use each for
the question it answers. A trace is not a replacement for a durable audit log,
and a metric cannot show every cause.

Every signal needs a schema and a cost model. Include request or trace identity,
route, outcome, dependency, and version where useful. Avoid high-cardinality
labels such as raw user IDs in metrics. Redact secrets and personal data before
emission, not only in a dashboard.

\section{Alerting}
An alert should indicate an actionable condition, have an owner, and explain the
first safe action. Alert on symptoms that matter to users or operators: an SLO
burn rate, queue age beyond a business limit, or inability to accept an order.
Avoid alerts that merely say a process is busy unless somebody can act on them.

A dashboard is a decision surface. Arrange it from outcome to cause: service
health, user journey, traffic, errors, latency, saturation, dependencies, and
recent changes. Test it during a failure; a beautiful dashboard that cannot
answer “what changed?” is not operationally useful.

\section{Runbooks and safe operations}
A runbook should say when it applies, what evidence to collect, what actions are
safe, what actions are irreversible, how to escalate, and how to verify recovery.
Prefer reversible controls: reduce traffic, disable a feature, pause a consumer,
or increase capacity. Administrative endpoints need authentication,
authorization, audit, and a break-glass procedure.

\section{Incident response}
During an incident, stabilize first. Establish a coordinator, a communications
channel, an incident timeline, and an explicit hypothesis. Separate observations
from guesses. Record commands and configuration changes. Avoid risky experiments
on a system whose state is not understood.

An incident is a period where the system's expected behaviour is not being met;
it is not limited to outages. A privacy exposure, incorrect financial result, or
unreliable deployment can be an incident even when servers respond normally.

\section{Postmortems}
A useful postmortem is blameless about individual intent and precise about system
conditions. Include impact, detection, timeline, contributing factors, what
worked, what failed, and actions with owners and dates. Do not produce only a
long list of “be more careful.” Improve a control: add a validation, limit blast
radius, make a signal actionable, rehearse recovery, or change an approval path.

\section{Change and observability}
Observability must evolve with the system. A new queue requires queue age and
dead-letter signals. A new authorization boundary requires denied-access audit
events. A new batch job needs freshness and completeness, not only process
uptime. Treat telemetry schemas as compatibility-sensitive APIs.

\begin{exercise}
Design an alert for an order service whose success SLO is burning budget rapidly.
Name the indicator, aggregation window, owner, first action, and one false
positive. Then write the minimum runbook needed to pause new intake safely.
\end{exercise}

\section{Chapter summary}
Operations is part of the product. Emit useful, privacy-aware signals; alert on
actionable user impact; maintain reversible runbooks; coordinate incidents with
evidence; and convert postmortem learning into stronger system controls.


\chapter{Refactoring, technical debt, legacy systems, and migrations}

\section{Maintenance is the normal state}
Most software cost occurs after the first release: bugs are fixed, dependencies
change, requirements grow, and people learn what the system really does. A design
that is easy to create but hard to change is not finished engineering.

\section{Refactoring}
Refactoring changes internal structure while preserving specified external
behaviour. Fowler's book popularized a catalog of small transformations and the
discipline of keeping the system working while its design improves
\cite{fowler2018refactoring}. The key precondition is a trustworthy safety net:
tests, assertions, domain invariants, or a controlled comparison.

A safe sequence is:

\begin{enumerate}
  \item characterize current behaviour, including surprising behaviour;
  \item make one structural change;
  \item run fast checks and inspect the diff;
  \item preserve or improve the boundary;
  \item repeat, then remove the obsolete path.
\end{enumerate}

Do not call a behaviour change a refactoring. A combined structural and product
change makes failures harder to attribute and rollback more difficult.

\section{Technical debt}
Technical debt is a metaphor for an intentional or accidental future cost. It
can be rational: a temporary adapter may enable a time-critical launch. It is
dangerous when the cost is not recorded, the interest is underestimated, or no
repayment trigger exists. Distinguish debt from mere code you dislike. A messy
but stable module may be cheaper than a fashionable rewrite.

Record debt with:

\begin{itemize}
  \item the shortcut and why it was taken;
  \item the risk or recurring cost;
  \item evidence of impact;
  \item the condition that makes repayment worthwhile;
  \item an owner or review date.
\end{itemize}

\section{Legacy systems}
Legacy means important software whose current behaviour is poorly understood or
whose constraints are hard to change. The first task is not to replace it. It
is to map dependencies, data, users, operational procedures, and hidden
contracts. Add characterization tests at the boundary. Capture current outputs
before changing internals.

A strangler approach can route one capability through a new implementation while
the old system remains. It reduces a big-bang cutover but creates two systems,
data synchronization, and ambiguous ownership during the transition.

\section{Migrations}
For a data migration, define source and target invariants, transformation rules,
handling of invalid rows, reconciliation totals, rollback or compensation, and
the cutover criterion. Run a dry run against representative data. Measure not
only whether rows copied but whether user-visible meaning was preserved.

For a service migration, use shadow traffic only when privacy and side effects
are controlled. Compare outputs with a defined equivalence relation. A new
implementation may be “more correct” and still break clients that depended on
old rounding, ordering, or error codes; decide which compatibility is required.

\section{When to rewrite}
A rewrite is justified by a concrete constraint: unsupported platform, unsafe
security posture, impossible performance target, or cost that exceeds a staged
transition. Estimate the work of reproducing undocumented behaviour, migrating
data, retraining operators, and running both systems. “The code is old” is not a
sufficient business case.

\begin{exercise}
Plan a migration from a legacy order table that stores money as floating point to
integer cents. Include discovery, rounding policy, dual read/write period,
reconciliation, rollback, and the point at which the old column can be removed.
\end{exercise}

\section{Chapter summary}
Maintenance is the main lifecycle, not a cleanup phase. Refactor in small safe
steps, record technical debt as a decision with a trigger, characterize legacy
behaviour, and migrate data with explicit invariants and reconciliation.


\chapter{Accessibility, internationalization, ethics, sustainability, and AI assistance}

\section{People are part of the operating environment}
Software affects people with different abilities, languages, devices, resources,
and power. Human requirements are not decorative polish. A service that cannot
be used with a keyboard, that assumes one date format, or that quietly denies a
group access has a functional defect for those users.

WCAG 2.2 organizes accessibility guidance around perceivable, operable,
understandable, and robust principles, with testable success criteria
\cite{w3c2023wcag}. Treat conformance as a useful baseline, not as proof that a
particular person can complete a task. Test with assistive technologies and with
people who use them when possible.

\section{Practical accessibility}
For a web interface, check semantic structure, heading order, labels, focus
visibility, keyboard operation, contrast, motion preferences, error recovery,
live-region announcements, and responsive zoom. Do not rely on color alone. Make
the primary task possible without a mouse and without a fast connection.

Accessibility must include operational interfaces. A dashboard with a tiny red-
green status indicator can exclude operators with color-vision differences and
can encourage unsafe decisions under pressure.

\section{Internationalization}
Separate translatable content from code, support plural rules, keep user locale
and time zone explicit, format numbers and dates according to context, and avoid
assuming names, addresses, or writing direction. Store canonical instants and
localize presentation. Test long strings, right-to-left layouts, non-ASCII input,
and culturally unfamiliar calendars where relevant.

\section{Sustainability}
Software consumes energy through computation, storage, networks, manufacturing,
and replacement cycles. A smaller carbon footprint can come from efficient
algorithms, right-sized infrastructure, caching, data retention limits, longer
hardware life, and avoiding work that has no user value. Optimization may trade
energy against latency or availability; measure the relevant system boundary
instead of claiming that one technique is always greener.

\section{Ethics and power}
Ask who benefits, who bears risk, who can contest a decision, and who can leave
the system. A model may be statistically accurate and still unfairly allocate
opportunity. A monitoring feature may improve security and also enable coercion.
Document purpose, limits, affected groups, appeal paths, and governance. Ethical
analysis is not solved by adding the word “responsible” to a project name.

\section{Responsible AI-assisted development}
AI assistants can accelerate drafting, search, and exploration. They can also
invent APIs, reproduce insecure code, leak confidential context, or create an
illusion of review. Treat generated code as an untrusted contribution:

\begin{enumerate}
  \item specify the behaviour and constraints before prompting;
  \item review generated code like an external patch;
  \item run tests, static analysis, dependency and secret checks;
  \item verify licenses, provenance, privacy, and security-sensitive logic;
  \item retain human ownership for design, approval, and incident response.
\end{enumerate}

Never accept a citation, benchmark, command, or library interface merely because
it sounds plausible. Check primary documentation and execute safe examples.
AI assistance is a tool for increasing feedback speed, not a transfer of
accountability.

\section{Inclusive design review}
Include affected perspectives before implementation, not only in a final audit.
Use a short risk register with severity, evidence, mitigation, owner, and review
date. Ask what happens to a person who has no high-end device, cannot hear an
alert, has a different name format, or is wrongly classified. These questions
often reveal ordinary usability and reliability improvements as well as ethical
risks.

\begin{exercise}
Choose a feature flag that automatically limits accounts after suspicious
activity. Identify accessibility, language, privacy, fairness, and appeal risks.
Design a human-review path and three metrics that could reveal harm.
\end{exercise}

\section{Chapter summary}
Design for real variation in people and contexts. Use accessibility standards as
a baseline, make locale assumptions explicit, measure sustainability trade-offs,
analyze power and contestability, and verify AI-assisted work with the same or
stronger discipline as human-authored code.


\chapter{Estimation, teams, incidents, and engineering management}

\section{Estimation under uncertainty}
An estimate is a prediction with assumptions, not a promise extracted from a
person. Separate three quantities: effort, elapsed time, and confidence. A task
may require two days of work but a week of elapsed time because it depends on an
approval or a test environment.

Use ranges when uncertainty is material. State the reference class, known work,
unknowns, dependencies, and what would change the estimate. Decompose until a
piece is small enough to observe progress, but do not confuse more detailed
fiction with more accuracy.

\section{Planning and sequencing}
Sequence work by risk and learning, not only by apparent feature value. Resolve
unknown integrations early. Build a vertical slice that exercises the real path
from input to persistence and observation. This uncovers hidden work sooner than
building isolated layers that cannot yet run.

A useful plan has:

\begin{itemize}
  \item outcome and non-goals;
  \item decision log and assumptions;
  \item milestones that produce observable evidence;
  \item dependencies and owners;
  \item quality and operational exit criteria;
  \item a change and rollback path.
\end{itemize}

\section{Team organization}
Team structure changes communication paths and ownership. Conway's law is a
useful observation about correspondence between communication structures and
system design, but it is not a deterministic theorem. Align ownership with
cohesive capabilities where possible. Make interfaces and escalation paths
explicit when ownership must cross teams.

Avoid “one person owns quality.” Every engineer affects quality, while specific
roles may coordinate standards, testing, security, or operations. A team is
responsible for the system it changes; a platform or specialist team should
reduce friction, not become a queue for every decision.

\section{Engineering management}
Management is system design applied to people, information, incentives, and
constraints. Good managers make priorities, decision rights, and trade-offs
clear; protect focus; develop capability; and ensure that uncomfortable evidence
can travel upward. Metrics should support learning rather than invite gaming.

Useful signals include lead time, change failure rate, recovery time, defect
escape, review wait, and user outcomes, but each can be misinterpreted. DORA-style
delivery measures are indicators, not a scoreboard for comparing unlike teams.
Combine quantitative signals with qualitative context.

\section{Incidents and learning}
The same skills used in operational incidents apply to project incidents:
surface facts, assign a coordinator, stabilize, communicate, and record decisions.
A missed migration, privacy concern, or repeated failed release deserves a causal
review. A postmortem action should improve the system, workflow, or decision
boundary; it should not merely ask people to remember harder.

\section{Decision rights}
For a recurring decision, name who recommends, who decides, who must be
consulted, and who is informed. Centralize decisions that need consistency or
rare expertise; delegate reversible local decisions. Ambiguous authority creates
delay and makes risk invisible.

\section{Healthy delivery pressure}
Urgency is real, but speed without a definition of done is borrowing from the
future. When a deadline forces a shortcut, record what is deferred, what risk it
creates, and the condition for repayment. A safe temporary reduction in scope is
usually easier to recover from than an invisible reduction in reliability or
security.

\begin{exercise}
Create a one-page plan for a four-person team delivering a new order export.
Include outcome, non-goals, risks, milestones, quality gates, owner boundaries,
and a rollback plan. Include one decision that should be delegated and one that
should be escalated.
\end{exercise}

\section{Chapter summary}
Estimate with ranges and assumptions. Sequence for learning and risk reduction,
align ownership with capabilities, use metrics as evidence rather than targets,
and turn incidents and shortcuts into explicit system improvements.


\chapter{Case study: engineering Ledgerline}

\section{Context}
Ledgerline accepts small orders from a web client and an internal support tool.
The first release has one team and modest traffic. Orders are commercially
important, must be auditable, and must not be duplicated by a retry. The team
chooses a modular monolith with a relational database. This is a deliberate
starting point, not a claim that every order system should remain a monolith.

\section{Requirements and boundaries}
The initial slice contains five obligations:

\begin{enumerate}
  \item validate positive quantities, positive integer prices, and a non-empty customer ID;
  \item calculate money exactly in integer cents;
  \item persist one accepted order and its items atomically;
  \item make retried creation idempotent through a unique key and stored result;
  \item expose health and operational signals without exposing personal data.
\end{enumerate}

The system explicitly does not include payment capture, inventory reservation,
or international tax law. Excluding those features is a safety decision: their
contracts and failure modes have not yet been designed.

\section{Architecture}
\begin{center}
\begin{tikzpicture}[node distance=4mm and 4mm,>=Latex,font=\small]
  \node[draw,rounded corners,fill=codebg,minimum width=21mm,minimum height=8mm] (client) {clients};
  \node[draw,rounded corners,fill=codebg,minimum width=25mm,minimum height=8mm,right=of client] (http) {HTTP adapter};
  \node[draw,rounded corners,fill=codebg,minimum width=29mm,minimum height=8mm,right=of http] (domain) {domain + application};
  \node[draw,rounded corners,fill=codebg,minimum width=22mm,minimum height=8mm,right=of domain] (db) {relational DB};
  \node[draw,rounded corners,fill=codebg,minimum width=27mm,minimum height=8mm,below=of domain] (telemetry) {telemetry + runbook};
  \draw[->,thick,accent] (client) -- (http);
  \draw[->,thick,accent] (http) -- (domain);
  \draw[->,thick,accent] (domain) -- (db);
  \draw[->,thick,accent] (http) |- (telemetry);
  \draw[->,thick,accent] (domain) -- (telemetry);
\end{tikzpicture}
\end{center}

The reference repository separates the domain calculation from the standard
library HTTP shell. The SQL schema repeats critical domain constraints. The
container runs as an unprivileged user with a read-only filesystem and a
healthcheck. These are modest controls, but each is visible and testable.

\section{API contract}
An illustrative request is:

\begin{lstlisting}[language={},caption={The body is a data contract, not a Python object.}]
POST /orders
Idempotency-Key: 8c4a...
Content-Type: application/json

{"customer_id":"customer-1",
 "items":[{"sku":"book","quantity":2,"unit_price_cents":1250}]}
\end{lstlisting}

On success, return a stable order identifier, status, and exact money values.
For a malformed request, return a client error without writing. For a reused
key with a different payload, reject the conflict rather than silently choosing
one request. For an unknown outcome after a timeout, let the caller retry with
the key and retrieve the stored result.

\section{Reviewing the implementation}
The Python domain module is short enough to inspect line by line:

\lstinputlisting[language=Python,caption={The domain calculation has explicit validation.}]{examples/python/ledgerline.py}

The code deliberately uses integer cents and a decimal library only for the
illustrative tax calculation. In a real financial system, tax policy, rounding
jurisdiction, currency, refunds, and audit requirements need a larger domain
model. The example is not a payment ledger.

The SQL schema is an independent guard. The test runner executes it in an
in-memory SQLite database and verifies that the view agrees with the inserted
line item. A production database choice changes locking, migrations, backups,
and operational procedures; the SQL illustrates principles rather than a vendor
recommendation.

\section{Failure analysis}
\noindent\begin{tabularx}{\textwidth}{@{}p{0.22\textwidth}p{0.27\textwidth}X@{}}
\toprule
Failure & Detection & Response \\
\midrule
Client retries after lost response & Idempotency-key lookup and duplicate metric & Return stored result; alert on unusual conflict rate. \\
Database unavailable & Error ratio, dependency latency, health check & Fail closed, preserve request identity, retry with bounded backoff. \\
Invalid quantity import & Database constraint and validation error & Reject transaction; record safe diagnostic. \\
Slow dependency added later & Trace and p95/p99 latency & Deadline, isolation, capacity review, or asynchronous workflow. \\
Secret in log & Secret scanning and log review & Revoke, scrub, investigate access, add regression check. \\
\bottomrule
\end{tabularx}

Each row names a detection and a response. “Add monitoring” would be too vague;
the signal has a source and an action.

\section{What would change at larger scale}
At higher traffic, the team might add a connection pool, read replicas, queue
workers, partitioning, or independent services. Each change creates new questions:
which data is authoritative, how are duplicates reconciled, how are versions
coordinated, and how are operators alerted? The existing boundary makes these
questions easier to isolate but does not answer them automatically.

\begin{exercise}
Extend Ledgerline with a payment authorization workflow. Add states for unknown
outcomes, a reconciliation job, an outbox event, and an audit record. Explain why
your design does not promise exactly-once effects across the payment provider.
\end{exercise}

\section{Chapter summary}
A small system can demonstrate serious engineering: explicit scope, exact money,
database invariants, idempotency, operational signals, unprivileged deployment,
and failure-oriented review. Scaling the system means revisiting constraints,
not replacing reasoning with more infrastructure.

\chapter{Capstone: from requirement to operated release}

\section{The assignment}
Design and implement a service that imports a customer order file, validates it,
stores accepted orders, exposes a query endpoint, and produces an operational
report. You may use the repository's Ledgerline examples as a starting point,
but the capstone should make its own decisions explicit.

\section{Required deliverables}
\begin{enumerate}
  \item a one-page problem statement with stakeholders, non-goals, and risks;
  \item an API and data specification with invalid examples;
  \item an architecture diagram and at least two ADRs;
  \item a working implementation with unit, component, and migration tests;
  \item a threat model and security verification checklist;
  \item a container image and CI workflow;
  \item an SLO, dashboard or query plan, and recovery runbook;
  \item a migration plan for one schema change;
  \item an incident scenario and blameless postmortem;
  \item a short reflection on accessibility, internationalization, sustainability,
        and AI-assisted development.
\end{enumerate}

\section{A release checklist}
\noindent\begin{longtable}{@{}p{0.07\textwidth}p{0.47\textwidth}p{0.28\textwidth}@{}}
\toprule
Done & Check & Evidence \\
\midrule
 & Behaviour and non-goals are clear & reviewed requirements \\
 & Invalid input is bounded and rejected & tests and validation code \\
 & Authorization is tested on allowed and denied paths & security tests \\
 & Data invariants exist at the right boundary & schema and migration checks \\
 & Retries and unknown outcomes are specified & contract and state model \\
 & Artifact is traceable to source and checks & CI record \\
 & Deployment is observable and reversible & runbook and rollout plan \\
 & Backups or exports can be restored & restore evidence \\
 & Accessibility and locale paths are tested & manual or automated evidence \\
 & Dependencies and licenses are reviewed & inventory and policy \\
\bottomrule
\end{longtable}

\section{How to present the design}
Begin with the user-visible outcome. Show one complete path, then one failure
path. State assumptions before diagrams. Explain why the design is sufficient
for the current constraints and what would force a different architecture.
Demonstrate a test failure or simulated incident rather than claiming that a
feature is “production ready” without evidence.

\section{Final perspective}
The central habit of software engineering is to turn invisible assumptions into
visible, testable, revisable decisions. Reliable systems are not produced by one
perfect design. They emerge from boundaries that fit the problem, feedback that
arrives before consequences become large, and teams that treat operation and
maintenance as part of building.

\begin{keyidea}
Build the smallest system that satisfies the important constraints, but make its
constraints explicit enough that a future team can safely discover when “smallest”
is no longer sufficient.
\end{keyidea}


\appendix
\chapter{Repository and tool guide}

\section{Quick start}
From the repository root:

\begin{lstlisting}[language=bash]
make test
make pdf
\end{lstlisting}

The first command runs standard-library checks. It compiles the C++ example with
warnings enabled, executes the TypeScript file with Node.js, validates SQL with
SQLite, and runs the shell check in a mocked mode. The second command runs
LaTeX, BibTeX, MakeIndex, and the final PDF build.

\section{Example inventory}
\noindent\begin{tabularx}{\textwidth}{@{}>{\raggedright\arraybackslash}p{0.3\textwidth}X@{}}
\toprule
Path & Purpose \\
\midrule
\path{examples/python/ledgerline.py} & Domain validation and exact money calculation \\
\path{examples/python/app.py} & Dependency-free health endpoint \\
\path{examples/typescript/order.ts} & Typed client-side calculation \\
\path{examples/java/RetryPolicy.java} & Bounded exponential backoff \\
\path{examples/cpp/latency_bucket.cpp} & Percentile selection over a sample \\
\path{examples/sql/schema.sql} & Relational constraints and an aggregate view \\
\path{examples/shell/healthcheck.sh} & Fail-fast health probe \\
\path{Dockerfile} and \path{compose.yaml} & Minimal container and local service orchestration \\
\bottomrule
\end{tabularx}

The examples are intentionally small. A small example can expose a rule without
requiring a framework, cloud account, or generated boilerplate. Production use
would require authentication, persistence lifecycle, request parsing, rate
limits, structured telemetry, deployment policy, and a threat model.

\section{Validation policy}
The repository distinguishes “verified here” from “instructions for another
environment.” The Python, SQL, shell, C++, and Node checks are run by the script
\path{scripts/verify_examples.py}. The Java source is dependency-free and includes a
compile command; this build environment has a Java runtime but not a compiler,
so the CI or a local JDK should compile it. The Docker setup is inspected as
source; Docker is optional in the current environment.

This distinction is important. A source file cannot be described as compiled
merely because it looks syntactically familiar. Record the toolchain and the
exact command used for verification.

\section{Updating the manuscript}
Edit chapter files rather than generated build files. Run `make pdf`, render the
PDF pages, inspect the title page, contents, code blocks, tables, and diagrams,
then run `make test`. If references change, rebuild BibTeX and the index. Keep
claims about versioned tools tied to a dated reference or a version range.

\chapter{Cumulative exercises}

\section{Foundations}
\begin{enumerate}
  \item Write a system boundary for a bicycle rental service. Name two actors outside the boundary and three invariants inside it.
  \item Convert three vague feature requests into requirements with measurable acceptance criteria.
  \item Draw a context diagram for a service that sends notifications through a third-party provider.
  \item Identify one stakeholder whose risk is not represented by the product requester.
\end{enumerate}

\section{Design and data}
\begin{enumerate}
  \setcounter{enumi}{4}
  \item Decompose a document-import service around changing decisions and state invariants.
  \item Design an API for cursor pagination. Explain cursor stability under inserts.
  \item Model a subscription lifecycle as a state machine and list illegal transitions.
  \item Write a schema constraint that prevents a child row from referring to a deleted parent.
  \item Design an expand-and-contract migration for renaming a public field.
\end{enumerate}

\section{Concurrency and architecture}
\begin{enumerate}
  \setcounter{enumi}{9}
  \item Find a lost-update interleaving in a stock reservation operation and repair it.
  \item Specify the retry and idempotency behaviour of a message consumer.
  \item Compare a modular monolith and two services for a team of three engineers.
  \item Write an ADR for choosing a relational database over a document store.
  \item State one architecture fitness function for each of three boundaries.
\end{enumerate}

\section{Quality and operation}
\begin{enumerate}
  \setcounter{enumi}{14}
  \item Build a test matrix for authorization and malformed input.
  \item Write a property that should hold for exact integer-money totals.
  \item Create an SLO with a defined measurement boundary and monthly error budget.
  \item Design a runbook for a growing queue and name the safe first action.
  \item Perform a postmortem for an accidentally deleted database column.
\end{enumerate}

\section{Security and human impact}
\begin{enumerate}
  \setcounter{enumi}{19}
  \item Threat-model a password-reset endpoint and include token replay.
  \item Identify personal data in logs, traces, backups, and analytics for an order system.
  \item Test a critical flow with keyboard-only navigation and a screen reader.
  \item List locale assumptions in a billing statement and replace them with explicit data.
  \item Review an AI-generated patch for invented APIs, insecure defaults, and license risk.
\end{enumerate}

\section{Capstone extensions}
\begin{enumerate}
  \setcounter{enumi}{24}
  \item Add an outbox and an idempotent event consumer to Ledgerline.
  \item Add a replayable migration with reconciliation and a failure-injection test.
  \item Produce a canary deployment plan with a rollback threshold.
  \item Measure the service under a stated workload and explain the bottleneck.
  \item Write a one-page architecture review that names what you would not build yet.
\end{enumerate}


\chapter{Selected solutions and reasoning}

\section{Foundations}
\begin{enumerate}
  \item A bicycle rental boundary can include reservation, payment authorization,
        and lock commands. The rider, payment provider, and physical lock are
        outside it. Invariants include one active rental per bicycle, a return
        cannot precede a start, and a charged amount matches the tariff snapshot.
  \item “Fast” becomes a journey, a percentile, a window, and a threshold. For
        example: p95 from submit to visible confirmation is below 400 ms for a
        defined payload at a stated load. Add invalid-input and dependency-failure
        acceptance cases.
  \item A notification diagram should show the caller, notification service,
        provider, delivery callback, and operator. The provider is a trust and
        reliability boundary; callbacks need authentication and replay handling.
\end{enumerate}

\section{Design and data}
\begin{enumerate}
  \setcounter{enumi}{4}
  \item Useful modules for document import are format parsing, validation,
        canonical domain model, persistence, and reporting. The parser should not
        decide authorization or write directly to the database.
  \item A cursor should encode or reference a stable ordering key and direction.
        It must not assume that page number remains stable while rows are inserted.
  \item A subscription state machine might allow trial -> active -> paused ->
        cancelled, but disallow cancelled -> active without a new subscription or
        explicit reactivation policy. The rule belongs in one transition function.
  \item A foreign key with `ON DELETE RESTRICT` prevents deletion while children
        exist; `ON DELETE CASCADE` is appropriate only when child loss is intended.
  \item Add the new field, write both fields, read the new field with fallback,
        backfill, enforce the new invariant, then stop writing and remove the old
        field after the compatibility window.
\end{enumerate}

\section{Concurrency and architecture}
\begin{enumerate}
  \setcounter{enumi}{9}
  \item Two reservations can both read stock 1. An atomic update such as
        `UPDATE stock SET available = available - 1 WHERE sku = ? AND available > 0`
        followed by checking affected rows prevents the lost update.
  \item A consumer acknowledges only after its effect is durable. If it crashes
        before acknowledgment, the message may repeat; use a unique event ID or
        idempotent upsert and accept that exactly-once effects are not implied.
  \item With three engineers and modest traffic, a modular monolith minimizes
        operational and coordination cost. A service split becomes justified by
        distinct ownership, legal isolation, or a measured scaling boundary.
  \item An ADR should name workload, transactions, operational skill, migration
        cost, and alternatives. “SQL is standard” is not enough evidence.
  \item A dependency rule, migration coexistence test, and tenant-isolation test
        are three possible fitness functions.
\end{enumerate}

\section{Quality and operation}
\begin{enumerate}
  \setcounter{enumi}{14}
  \item Cover authenticated allowed access, authenticated denied access,
        unauthenticated access, malformed identifiers, oversized input, and
        dependency failure. State which combinations are intentionally omitted.
  \item For every valid item list, total equals the sum of positive integer
        quantity-times-price products and remains within the chosen numeric bound.
  \item An SLO must identify eligible requests, success definition, latency
        boundary, window, and excluded maintenance if any. Its budget is the
        complement of the target over the eligible measurement population.
  \item First stop or rate-limit intake if backlog threatens data freshness, then
        inspect consumer errors, dependency latency, and poison messages before
        increasing concurrency.
  \item A postmortem should preserve evidence, identify the deletion path and
        missing guard, restore from a verified backup, and add a migration check
        plus permission or review control. Blaming the operator is not a control.
\end{enumerate}

\section{Security and human impact}
\begin{enumerate}
  \setcounter{enumi}{19}
  \item Use short-lived, single-purpose reset tokens stored as hashes, bind them
        to the account and intent, invalidate on use, rate-limit attempts, avoid
        account enumeration, and log safe audit events.
  \item Remove secrets and unnecessary personal data from telemetry; define access,
        retention, deletion, and break-glass review for remaining records.
  \item Check semantic controls, focus order, visible focus, labels, error
        announcements, and completion without a pointer device. Confirm with
        assistive technology rather than relying on a visual inspection.
  \item Store canonical instants, currency, locale, and user preferences
        separately. Localize format at the presentation boundary.
  \item Verify every claimed API and dependency, run tests and security checks,
        inspect data handling, and require a human owner for approval and release.
\end{enumerate}

\section{Capstone reasoning}
\begin{enumerate}
  \setcounter{enumi}{24}
  \item An outbox row and the order state belong in one local transaction. The
        consumer uses a unique event ID and records completion. Re-delivery is
        expected; the external effect must be idempotent or reconciled.
  \item A migration should be rerunnable or safely resumable, report rejected
        rows, compare source and target totals, and stop before destructive work
        if reconciliation fails.
  \item A canary needs a cohort, success indicators, a time window, an owner,
        and an automatic or human-controlled rollback action. Do not expose a
        wider audience before telemetry is healthy.
  \item Measure a realistic workload and identify the saturated resource. A
        higher average throughput number without tail and error data is incomplete.
  \item The best architecture review names deferred features and the signals that
        would justify them. Restraint is part of design.
\end{enumerate}



\backmatter
\chapter*{Glossary}
\addcontentsline{toc}{chapter}{Glossary}
\footnotesize
\begin{description}[style=nextline,leftmargin=1.35in,labelwidth=1.2in]
\item[API] A defined interface through which software components exchange requests and results.
\item[Availability] The proportion of time a service is usable under a stated measurement definition.
\item[Capacity] The sustainable amount of work a system can handle while meeting its objectives.
\item[Coupling] The degree to which one component depends on another component's details or timing.
\item[Deployment] The act of making a version available in an execution environment.
\item[Idempotency] The property that repeating an operation has the same intended effect as performing it once.
\item[Invariant] A condition that remains true across a specified operation or state transition.
\item[Latency] The elapsed time between an operation's start and completion, measured at a defined boundary.
\item[Observability] The ability to infer internal system state from externally emitted signals.
\item[Recovery point objective] The maximum acceptable amount of data loss measured in time.
\item[Recovery time objective] The maximum acceptable time to restore a service after a failure.
\item[Reliability] The probability that a system performs its required function for a stated period and context.
\item[Service-level objective] A target for a measurable service indicator over a stated window.
\item[Technical debt] A present design or implementation choice whose future cost is knowingly incurred to gain a short-term benefit.
\item[Threat model] A structured description of assets, actors, trust boundaries, threats, and mitigations.
\end{description}

\bibliographystyle{plain}
\bibliography{references}

\printindex
\end{document}
